% ============================================================
%  INSUMO DE TESIS — Metodología, Código y Cifras
%  Pensión 65 y Adopción de Billetera Digital (ENAHO 2024)
%  Autor: Cecilia Vargas Risco
%  Generado: 2026-07-02
% ============================================================
\documentclass[12pt, a4paper]{article}
\usepackage[utf8]{inputenc}
\usepackage[T1]{fontenc}
\usepackage[spanish]{babel}
\usepackage{geometry}
\geometry{margin=2.5cm}
\usepackage{booktabs}
\usepackage{longtable}
\usepackage{array}
\usepackage{multirow}
\usepackage{amsmath}
\usepackage{amssymb}
\usepackage{xcolor}
\usepackage{hyperref}
\usepackage{graphicx}
\usepackage{float}
\usepackage{enumitem}
\usepackage{fancyhdr}
\usepackage{titlesec}
\usepackage{lmodern}
\usepackage{microtype}
\usepackage{parskip}
\usepackage{caption}
\usepackage{subcaption}
\usepackage{threeparttable}
\usepackage{pdflscape}
\usepackage{afterpage}
\usepackage{tcolorbox}
\tcbuselibrary{skins, breakable}

\hypersetup{
  colorlinks=true,
  linkcolor=blue!60!black,
  urlcolor=blue!60!black,
  citecolor=green!50!black
}

\pagestyle{fancy}
\fancyhf{}
\rhead{\small Insumo de Tesis — Pensión 65 y Billetera Digital}
\lhead{\small Vargas Risco, 2026}
\cfoot{\thepage}

\titleformat{\section}{\large\bfseries}{}{0em}{}[\titlerule]
\titleformat{\subsection}{\normalsize\bfseries}{}{0em}{}

% Cajas de interpretación
\newtcolorbox{interpcaja}[1][]{
  colback=blue!5!white,
  colframe=blue!40!black,
  title=#1,
  fonttitle=\bfseries\small,
  breakable
}
\newtcolorbox{alertcaja}[1][]{
  colback=orange!8!white,
  colframe=orange!70!black,
  title=#1,
  fonttitle=\bfseries\small,
  breakable
}
\newtcolorbox{resultcaja}[1][]{
  colback=green!5!white,
  colframe=green!50!black,
  title=#1,
  fonttitle=\bfseries\small,
  breakable
}

% ============================================================
\begin{document}
% ============================================================

\begin{titlepage}
\centering
\vspace*{2cm}
{\LARGE\bfseries Insumo de Tesis\\[0.4em]}
{\Large Pensión 65 y Adopción de Billetera Digital\\[0.3em]}
{\large Metodología, Código y Cifras Verificadas\\[2cm]}
{\normalsize Cecilia Vargas Risco\\[0.5em]}
{\small svargasrisco@gmail.com\\[0.5em]}
{\small Julio 2026\\[3cm]}
\begin{tcolorbox}[colback=gray!10, colframe=gray!40, width=0.85\textwidth]
\small\centering
Este documento consolida toda la metodología, las variables construidas,
el pipeline de estimación y los resultados numéricos del análisis empírico
sobre Pensión 65 y adopción de billetera digital usando ENAHO 2024.
Está pensado como insumo directo para redactar las secciones 3--6 del documento de tesis.
\end{tcolorbox}
\end{titlepage}

\tableofcontents
\newpage

% ============================================================
\section{Pregunta de Investigación y Diseño General}
% ============================================================

\subsection{Pregunta central}

¿Recibir Pensión 65 ---el programa de transferencia monetaria no condicionada
del Estado peruano para adultos mayores en situación de pobreza extrema---
aumenta la probabilidad de adoptar una billetera digital (Yape, Plin o Cuenta
DNI) entre los beneficiarios elegibles?

\subsection{Estrategia de identificación}

El programa Pensión 65 define la elegibilidad mediante dos condiciones
simultáneas: (1)~tener 65 años o más, y (2)~estar clasificado como pobre
extremo por el Sistema de Focalización de Hogares (SISFOH) del Ministerio
de Desarrollo e Inclusión Social (MIDIS). La condición de edad crea una
\textbf{discontinuidad exógena}: individuos justo por encima y justo por
debajo de los 65 años son comparables en características observables y no
observables, pero solo los primeros son elegibles para el programa.

El diseño empleado es un \textbf{Diseño de Regresión en Discontinuidad (RDD)},
tomando la edad en años como variable de corte (\textit{running variable})
y el umbral de 65 años como punto de quiebre. En su versión \textit{fuzzy},
el corte de edad actúa como instrumento para la recepción efectiva del bono,
estimando el efecto local promedio del tratamiento (LATE) sobre los
\textit{compliers} ---aquellos que reciben el programa gracias a haber
cumplido 65 años.

\begin{alertcaja}[Error conceptual corregido en este análisis]
El paper original restringía la muestra a \texttt{POBREZA=1} (pobreza
monetaria extrema según la ENAHO, calculada por el INEI en función del gasto
per cápita frente a la línea alimentaria) y la etiquetaba como ``SISFOH
extrema pobreza''. Esto es incorrecto: POBREZA es una clasificación
\textit{ex post} basada en gasto corriente, mientras que SISFOH es un
\textit{proxy means test} basado en características predeterminadas de
vivienda. De los 4{,}437 receptores reales de Pensión 65 (con 65+ años)
en ENAHO 2024, solo el 11.3\% tiene \texttt{POBREZA=1}. El 88.7\% restante
tiene mayor gasto porque el bono mismo eleva el consumo del hogar
(causalidad reversa). Por tanto, el análisis primario usa la muestra
completa (\textit{full sample}) sin restricción por \texttt{POBREZA}.
\end{alertcaja}

% ============================================================
\section{Fuente de Datos: ENAHO 2024}
% ============================================================

\subsection{Encuesta Nacional de Hogares 2024}

La fuente primaria es la Encuesta Nacional de Hogares sobre Condiciones de
Vida y Pobreza (ENAHO), edición 2024, producida por el Instituto Nacional
de Estadística e Informática (INEI). Es la encuesta de hogares más completa
del Perú: representativa a nivel nacional, con panel rotativo y diseño
probabilístico estratificado.

\subsection{Módulos utilizados}

Se emplean seis módulos de la ENAHO 2024:

\begin{longtable}{@{}llp{7cm}l@{}}
\toprule
Módulo & Nombre & Contenido relevante & Filas \\\midrule
\endhead
01 & Vivienda & Paredes, piso, techo, agua, desagüe, combustible, electricidad, internet & 44{,}731 hogares \\
02 & Características de los miembros & Edad, sexo, área, factor de expansión (\texttt{FACPOB07}) & 117{,}721 personas \\
03 & Educación & Nivel educativo del jefe y todos los miembros (\texttt{P301A}) & 106{,}619 personas \\
05 & Empleo e ingresos & Billetera digital (\texttt{P558E1\_9}, \texttt{P558Hk\_7}), Pensión 65 (\texttt{P5567A}), ocupación, bancarización & 85{,}992 personas \\
18 & Equipamiento del hogar & Tenencia de smartphone, TV, refrigeradora, lavadora, computadora & 943{,}348 registros (formato largo) \\
34 & Sumaria & Ingresos del hogar, transferencias (\texttt{INGTPU03}, \texttt{INGTPU01}), clasificación POBREZA, factor de expansión & 33{,}691 hogares \\
\bottomrule
\caption{Módulos ENAHO 2024 utilizados en el análisis.}
\end{longtable}

\subsection{Llaves de merge}

El merge entre módulos se realiza por:
\begin{itemize}[noitemsep]
  \item Nivel hogar: \texttt{CONGLOME + VIVIENDA + HOGAR}
  \item Nivel persona: \texttt{CONGLOME + VIVIENDA + HOGAR + CODPERSO}
\end{itemize}

El dataset analítico final resulta de cruzar los módulos de vivienda, educación
y Sumaria (nivel hogar) con el módulo 05 (nivel persona), priorizando la
dimensión persona como unidad de análisis. El módulo 18 se pivota de formato
largo a ancho antes del merge.

\subsection{Muestra resultante}

\begin{center}
\begin{tabular}{lrr}
\toprule
Etapa & Observaciones & Descripción \\\midrule
Raw Módulo 02 & 117{,}721 & Todos los miembros de hogar \\
Después de merge con M05 & 85{,}992 & Solo personas con información de empleo \\
Dataset analítico final (M01+02+03+05+18+34) & 117{,}721 & Merge completo \\
Después de filtrar EDAD no nula & 113{,}755 & Muestra RDD base \\
\quad Grupo de control (EDAD $<$ 65) & 99{,}667 & \\
\quad Grupo de tratamiento (EDAD $\geq$ 65) & 14{,}088 & \\
\bottomrule
\end{tabular}
\end{center}

% ============================================================
\section{Variables Construidas}
% ============================================================

\subsection{Variable dependiente: TIENE\_BILLETERA}

\textbf{Definición:} indicador binario que vale 1 si la persona tiene o usa
activamente una billetera digital (Yape, Plin o Cuenta DNI).

\textbf{Construcción:} combinación OR de dos componentes del Módulo 05:

\begin{enumerate}[noitemsep]
  \item \textbf{Tenencia} (\texttt{P558E1\_9 == ``9''}): la pregunta P558E1
    lista los instrumentos financieros que el encuestado posee; la categoría
    9 corresponde a ``billetera digital''.
  \item \textbf{Uso activo} (\texttt{P558Hk\_7 == ``7''} para
    $k = 1, \ldots, 12$): la pregunta P558Hk pregunta con qué tipo de
    instrumento financiero realizó pagos en el mes $k$ anterior; la
    categoría 7 corresponde a ``billetera digital''.
\end{enumerate}

$$\texttt{TIENE\_BILLETERA} = (\texttt{tenencia} = 1) \vee (\exists k \in [1,12] : \texttt{P558Hk\_7} = 7)$$

\textbf{Variable complementaria:} \texttt{USA\_BILLETERA} captura solo el
componente de uso activo (sin tenencia pasiva), siendo un indicador más
restrictivo de adopción.

\textbf{Estadísticas descriptivas:}

\begin{center}
\begin{tabular}{lcc}
\toprule
Indicador & Porcentaje & N \\\midrule
Tenencia (P558E1\_9) & 10.1\% & 113{,}755 \\
Uso activo (P558Hk\_7, algún mes) & 21.3\% & 113{,}755 \\
\textbf{TIENE\_BILLETERA (combinado)} & \textbf{27.0\%} & \textbf{113{,}755} \\
\midrule
Grupo control (EDAD $<$ 65) & 22.2\% & 99{,}667 \\
Grupo tratado (EDAD $\geq$ 65) & 7.9\% & 14{,}088 \\
\bottomrule
\end{tabular}
\end{center}

\begin{interpcaja}[Interpretación para la tesis]
La tasa de adopción combinada del 27\% es consistente con la estadística
oficial del BCRP (2024) que reporta~46\% de adopción nacional para la
población general. La brecha entre control (22\%) y tratados (7.9\%) refleja
el diferencial estructural por edad: los adultos mayores tienen menor
familiaridad con tecnología digital, menor acceso a smartphones y menor
escolaridad. Esta correlación cruda negativa (mayores adoptan menos) no es
el efecto causal del programa, sino el efecto edad. El RDD precisamente
separa ambos efectos comparando individuos de edades muy similares alrededor
de los 65 años.
\end{interpcaja}

\subsection{Variable de tratamiento endógeno: RECIBE\_P65\_PERSONA}

\textbf{Código ENAHO:} \texttt{P5567A} del Módulo 05.

\textbf{Pregunta:} El ítem 7 del bloque P556iA pregunta si la persona recibe
el bono de Pensión 65. La respuesta 1 indica recepción.

\textbf{Construcción:}
$$\texttt{RECIBE\_P65\_PERSONA} = \mathbf{1}(\texttt{P5567A} = 1)$$

\textbf{Validación:}
\begin{itemize}[noitemsep]
  \item Receptores en muestra: \textbf{4{,}437 personas}
  \item A nivel hogar (INGTPU03 $>$ 0 en Sumaria): \textbf{3{,}564 hogares}
  \item Monto mediano del bono (INGTPU03): S/ 1{,}506/año $\approx$ S/ 251/mes
    $\approx$ S/ 250 del diseño oficial del programa \checkmark
  \item Expansión poblacional (con FACPOB07): $\approx$ 902{,}000 receptores,
    consistente con las 750{,}000 pensiones activas reportadas por MIDIS
\end{itemize}

\textbf{Correlación con la running variable:} la correlación de
\texttt{RECIBE\_P65\_PERSONA} con \texttt{EDAD} es $|r| = 0.403$, la más
alta entre todas las variables analíticas después de \texttt{CODPERSO}
(identificador que correlaciona por diseño muestral). Esto confirma que la
recepción del bono aumenta fuertemente con la edad, como se espera.

\textbf{Primera etapa (first stage):}

\begin{center}
\begin{tabular}{crr}
\toprule
Edad & N en muestra & \% recibe Pensión 65 \\\midrule
63 & 1{,}032 & 0.0\% \\
64 & 1{,}070 & 0.0\% \\
\midrule
\textbf{65} & \textbf{1{,}023} & \textbf{3.8\%} \\
66 & 887 & 16.0\% \\
70 & 805 & 26.6\% \\
80 & 379 & 45.9\% \\
\bottomrule
\end{tabular}
\end{center}

El salto de 0\% a 3.8\% exactamente en los 65 y de 3.8\% a 16.0\% en los
66 refleja el proceso de incorporación gradual al programa. El diseño fuzzy
capta este salto como primera etapa del instrumento de edad.

\subsection{Variables de control (covariables del RDD)}

\begin{longtable}{@{}llp{6.5cm}l@{}}
\toprule
Variable & Tipo & Descripción & Fuente \\\midrule
\endhead
\texttt{INTERNET\_HOGAR} & Binaria & El hogar tiene acceso a internet & M01 \\
\texttt{SMARTPHONE} & Binaria & El hogar posee al menos un smartphone & M18 \\
\texttt{INGRESO\_PC} & Continua & Ingreso total del hogar / nro. de miembros (S/ anuales) & Sumaria \\
\texttt{NIVEL\_EDUCATIVO} & Ordinal (1--12) & Nivel educativo del encuestado & M03 \\
\bottomrule
\caption{Covariables del RDD. Nota: \texttt{POBREZA} fue eliminada de las covariables porque es post-tratamiento (el bono eleva el gasto del hogar, lo que puede cambiar la clasificación de pobreza).}
\end{longtable}

\textbf{Estadísticas descriptivas por grupo:}

\begin{center}
\begin{tabular}{lcccc}
\toprule
Variable & \multicolumn{2}{c}{Tratados ($\geq$65)} & \multicolumn{2}{c}{Control ($<$65)} \\
 & Media & SD & Media & SD \\\midrule
TIENE\_BILLETERA & 0.079 & 0.269 & 0.222 & 0.415 \\
USA\_BILLETERA & 0.032 & 0.176 & 0.179 & 0.384 \\
INTERNET\_HOGAR & 0.417 & 0.493 & 0.464 & 0.499 \\
SMARTPHONE & 0.831 & 0.375 & 0.870 & 0.337 \\
INGRESO\_PC (S/ anuales) & 13{,}454 & 13{,}063 & 11{,}417 & 11{,}201 \\
NIVEL\_EDUCATIVO & 4.39 & 3.43 & 5.45 & 3.90 \\
N & \multicolumn{2}{c}{14{,}088} & \multicolumn{2}{c}{99{,}667} \\
\bottomrule
\end{tabular}
\end{center}

\subsection{Variable de transferencia placebo: RECIBE\_JUNTOS\_HOGAR}

\texttt{INGTPU01} de Sumaria identifica hogares receptores del programa
Juntos (transferencia condicionada a familias con niños). Esta variable es
un \textbf{placebo}: Juntos no tiene corte de elegibilidad a los 65 años,
por lo que no debería mostrar discontinuidad en el cutoff.

Receptores en muestra: \textbf{4{,}129 hogares}. Esta variable se usa en
las pruebas de robustez para verificar que el cutoff de edad no genera
discontinuidades espurias en transferencias no relacionadas.

\subsection{Índice de Focalización de Hogares (IFH)}

El IFH fue construido replicando la metodología de Bernal, Carpio y Klein
(2017, \textit{Journal of Public Economics}, Apéndice F, Tablas A.8/A.9/A.10).
Es el puntaje que usa el SISFOH-2010 para clasificar a los hogares como
elegibles para programas sociales.

\textbf{Estructura:} suma ponderada de variables categóricas de vivienda,
servicios básicos y activos del hogar, diferenciada por área geográfica:

\begin{center}
\begin{tabular}{llp{7cm}}
\toprule
Área & Nro. variables & Variables incluidas \\\midrule
Lima Metropolitana & 11 & Combustible, agua, paredes, techo, piso, desagüe, seguro médico, bienes, educación jefe, hacinamiento, teléfono fijo \\
Otras áreas urbanas & 10 & Igual a Lima sin teléfono fijo \\
Áreas rurales & 7 & Combustible, seguro médico, bienes, educación jefe, educación máxima del hogar, electricidad, piso tierra \\
\bottomrule
\end{tabular}
\end{center}

\textbf{Estandarización:} el IFH\_RAW (suma ponderada) se estandariza en el
rango [0, 100] por departamento (25 departamentos como aproximación de los
15 clusters SISFOH-2010).

\textbf{Umbral de elegibilidad:} varía por cluster/departamento según la
Tabla A.10 de BCK (2017), con rango entre 33 y 55 puntos. Promedio en
ENAHO 2024: umbral $= 45.68$ (rango: 33--55).

\textbf{Criterio SISFOH completo:}
$$\texttt{ELEGIBLE\_SIS} = (\texttt{IFH} \leq \texttt{UMBRAL}) \wedge (\text{gasto agua} \leq 20) \wedge (\text{gasto electricidad} \leq 25)$$

\textbf{Nota crítica:} al construir la variable de miembros con seguro médico,
se excluye el SIS (P4195) para evitar endogeneidad circular: si se incluyera,
la variable IFH dependería de la propia cobertura SIS que el programa
dispara (BCK 2017, footnote 2).

\textbf{Estadísticas del IFH en adultos mayores ($\geq$ 65 años):}
\begin{center}
\begin{tabular}{lr}
\toprule
Estadístico & Valor \\\midrule
N con IFH no nulo & 13{,}766 \\
Media del IFH & 60.52 \\
Desviación estándar & 19.07 \\
Mínimo & 0.00 \\
Máximo & 100.00 \\
Umbral promedio (UMBRAL) & 45.68 \\
Elegibles IFH (IFH $\leq$ UMBRAL) & 3{,}377 \\
Elegibles SIS (IFH + agua + electr.) & 2{,}042 \\
\bottomrule
\end{tabular}
\end{center}

\subsection{Muestras de análisis (SAMPLE\_FLAG)}

El pipeline construye cuatro submuestras identificadas por banderas:

\begin{center}
\begin{tabular}{llrr}
\toprule
Flag & Criterio & N & Receptores P65 (\%) \\\midrule
A & Muestra completa (primaria) & 113{,}755 & 4{,}437 (100\%) \\
B & POBREZA = 1 (extrema pobreza monetaria) & 6{,}717 & 501 (11.3\%) \\
C & POBREZA $\in \{1, 2\}$ (pobres monetarios) & 28{,}737 & 1{,}619 (36.5\%) \\
D & ELEGIBLE\_SIS = 1 (elegibles IFH BCK 2017) & 12{,}116 & 1{,}171 (26.4\%) \\
\bottomrule
\end{tabular}
\end{center}

\begin{interpcaja}[Interpretación para la tesis]
La Muestra A es el estimando primario: estima el efecto del programa en la
población general que rodea el cutoff de edad 65. Las muestras B y C son
\textbf{análisis de heterogeneidad descriptivos}, no diseños causales
primarios, porque POBREZA es una variable post-tratamiento (endógena al
bono). La Muestra D aproxima la elegibilidad SISFOH real usando el IFH
de BCK (2017) y tiene validez causal si el IFH es una buena proxy del
puntaje SISFOH oficial.
\end{interpcaja}

% ============================================================
\section{Pipeline del Código}
% ============================================================

El análisis está implementado en un único script Python modular:
\texttt{scripts/script\_completo/script.py}. Se estructura en 6 fases más
un bloque auxiliar:

\begin{center}
\begin{tabular}{lll}
\toprule
Fase & Función & Output principal \\\midrule
0 & \texttt{phase\_0\_download\_from\_inei()} & CSVs crudos descargados desde INEI \\
1 & \texttt{phase\_1\_preprocess()} & \texttt{enaho\_2024\_clean.csv} (117{,}721 × 47) \\
2 & \texttt{phase\_2\_clean()} & \texttt{enaho\_rdd\_full.csv} (113{,}755 × 57) \\
3 & \texttt{phase\_3\_main()} & \texttt{main\_results.csv} (18 estimaciones) \\
3b & \texttt{phase\_3b\_bloque\_f()} & \texttt{rdd2\_results.csv} (7 estimaciones) \\
4 & \texttt{phase\_4\_robustness()} & \texttt{robustness\_results.csv} (21 tests) \\
5 & \texttt{phase\_5\_output()} & Tablas LaTeX + figuras PDF \\
\bottomrule
\end{tabular}
\end{center}

\subsection{Fase 1 — Preprocesamiento}

Carga y consolida los 6 módulos. Los pasos clave son:
\begin{enumerate}[noitemsep]
  \item Cargar M01 (vivienda), M02 (miembros), M03 (educación), M05
    (empleo/billetera), M18 (equipamiento), Sumaria.
  \item Construir \texttt{TIENE\_BILLETERA} y \texttt{USA\_BILLETERA} desde M05.
  \item Extraer \texttt{RECIBE\_P65\_PERSONA} (\texttt{P5567A == 1}).
  \item Pivotar M18 de formato largo a ancho para obtener smartphone, TV, etc.
  \item Calcular \texttt{HACINAMIENTO = MIEPERHO / P104} (miembros / cuartos).
  \item Extraer \texttt{INGTPU03} y \texttt{INGTPU01} desde Sumaria.
  \item Hacer merge final a nivel persona usando las llaves de hogar + \texttt{CODPERSO}.
\end{enumerate}

\subsection{Fase 2 — Limpieza RDD}

\begin{enumerate}[noitemsep]
  \item Filtrar a observaciones con EDAD no nula: 117{,}721 $\to$ 113{,}755.
  \item Centrar la running variable: \texttt{RUNNING\_VAR = EDAD - 65}.
  \item Cargar \texttt{ifh\_2024.csv} (117{,}721 filas sin llaves de merge)
    y hacer un join posicional con las llaves de \texttt{enaho\_2024\_clean.csv}
    (mismo orden de filas, generados en una sola pasada).
  \item Crear las cuatro \texttt{SAMPLE\_FLAG}.
  \item Guardar \texttt{enaho\_rdd\_full.csv} con 57 columnas.
\end{enumerate}

\subsection{Fase 3 — Estimación principal}

Corre cinco bloques de estimación. La función central es \texttt{run\_rdd()}:
\begin{itemize}[noitemsep]
  \item Usa \texttt{rdrobust} con kernel triangular y bandwidth MSE-óptimo
    (\texttt{bwselect = ``mserd''}).
  \item Acepta un argumento \texttt{fuzzy\_col} para el diseño Fuzzy RDD.
  \item Si \texttt{rdrobust} falla (por matrices de tipo incompatible), cae
    en un estimador WLS local-lineal de statsmodels como fallback.
  \item Devuelve siempre un dict completo: estimador, SE robusto, IC 95\%,
    BW, N efectivo, método.
\end{itemize}

\subsection{Fase 4 — Robustez}

Corre 21 tests de sensibilidad: McCrary/rddensity, placebos de covariables,
sensibilidad al BW (×0.5, ×1, ×2), donut holes ($\pm$0.5, $\pm$1, $\pm$2
años), heterogeneidad por POBREZA, balance de covariables, sensibilidad al
polinomio (lineal vs. cuadrático), placebos de cutoff, e inferencia por
permutación (500 iteraciones).

\subsection{Bloque F — RDD2 con IFH como running variable}

Implementa el diseño alternativo de Bando, Galiani y Gertler (2020): en
vez de usar la edad como discontinuidad, usa el \textbf{IFH centrado en el
umbral departamental} ($\texttt{IFH} - \texttt{UMBRAL}$) como running
variable. Los hogares justo por debajo del umbral son elegibles SISFOH;
los que están justo por encima no lo son. La muestra se restringe a
adultos mayores ($\geq$ 65 años) para asegurar que la condición de edad
ya se cumpla.

% ============================================================
\section{Diseños de Estimación}
% ============================================================

\subsection{Diseño A — Sharp RDD (ITT), muestra completa}

\textbf{Especificación econométrica:}

$$Y_i = \alpha + \tau \cdot \mathbf{1}(\text{EDAD}_i \geq 65) + \beta_1 (\text{EDAD}_i - 65) + \beta_2 \mathbf{1}(\geq 65) \cdot (\text{EDAD}_i - 65) + \varepsilon_i$$

donde $Y_i \in \{\texttt{TIENE\_BILLETERA}_i, \texttt{USA\_BILLETERA}_i\}$.

\begin{itemize}[noitemsep]
  \item \textbf{Estimando:} Intención de Tratamiento (ITT) --- efecto de
    \textit{ser elegible} (cumplir 65 años) sobre la adopción, promediado
    sobre toda la población cercana al corte.
  \item \textbf{Kernel:} triangular (mayor peso a observaciones cercanas al cutoff).
  \item \textbf{Bandwidth:} MSE-óptimo, $h^* = 18.04$ años para TIENE\_BILLETERA.
  \item \textbf{SE:} HC1 robustos con sesgo corregido (rdrobust).
  \item \textbf{Especificaciones:} baseline (sin covariables), covariates
    (con INTERNET\_HOGAR, SMARTPHONE, INGRESO\_PC, NIVEL\_EDUCATIVO), y
    dpto\_fe (con efectos fijos de departamento vía WLS local).
\end{itemize}

\subsection{Diseño B — Fuzzy RDD (LATE), muestra completa}

\textbf{Especificación:} RDD con variable endógena.

\textbf{Primera etapa:}
$$\texttt{RECIBE\_P65\_PERSONA}_i = \pi_0 + \pi_1 \cdot \mathbf{1}(\text{EDAD}_i \geq 65) + f(\text{EDAD}_i - 65) + u_i$$

\textbf{Segunda etapa (Wald/IV):}
$$\tau_{\text{LATE}} = \frac{\tau_{\text{RF}}}{\pi_1} = \frac{\text{Forma reducida}}{\text{Primera etapa}}$$

\begin{itemize}[noitemsep]
  \item \textbf{Estimando:} Efecto Local Promedio del Tratamiento (LATE) ---
    efecto causal de \textit{recibir} el bono sobre los \textit{compliers}
    (quienes reciben porque cumplieron 65 años).
  \item \textbf{Implementación:} \texttt{rdrobust(y, x, fuzzy=RECIBE\_P65\_PERSONA)}.
  \item \textbf{Bandwidth:} $h^* = 6.45$ años (menor que el sharp RDD porque
    estima dos ecuaciones simultáneamente).
\end{itemize}

\begin{alertcaja}[Advertencia: poder estadístico del Fuzzy RDD]
El BW cae de 18 a 6.4 años y el N\_eff de 35{,}006 a 12{,}616. En ese rango
estrecho la primera etapa es débil (pocos compliers exactamente en el
cutoff), lo que produce un denominador pequeño en el Wald y SE muy grandes
($\approx 0.30$). Esto no es un error del código sino una característica
del dato: en la \textit{full sample} hay muy poca recepción del bono en
$[63, 67]$ años. Un fuzzy RDD con primera etapa fuerte requeriría una
submuestra más restringida (por ejemplo, Muestra D).
\end{alertcaja}

\subsection{Diseño C — RDD en muestra ELEGIBLE\_SIS=1 (Bloque F\_SAMPLE)}

Mismo sharp RDD que el Diseño A, pero restringido a las 12{,}116 personas
con \texttt{ELEGIBLE\_SIS = 1} (IFH $\leq$ UMBRAL según BCK 2017). Aproxima
el estimando del programa restringiendo la muestra a los hogares que el
SISFOH clasificaría como elegibles.

\subsection{Diseño D — RDD2: IFH como running variable (Bloque F)}

Replica la estrategia de identificación de Bando, Galiani y Gertler (2020)
para el contexto peruano:

$$Y_i = \alpha + \tau \cdot \mathbf{1}(\texttt{IFH}_i \leq \texttt{UMBRAL}_i) + g(\texttt{IFH}_i - \texttt{UMBRAL}_i) + \varepsilon_i$$

\begin{itemize}[noitemsep]
  \item \textbf{Running variable:} IFH $-$ UMBRAL (centrada en 0 en el umbral
    departamental).
  \item \textbf{Cutoff:} 0 (IFH exactamente igual al umbral).
  \item \textbf{Muestra:} adultos mayores ($\geq$ 65 años) con IFH no nulo
    ($N = 13{,}766$).
  \item \textbf{Estimando:} efecto de ser clasificado como elegible SISFOH
    (cruzar el umbral IFH) sobre la adopción de billetera digital.
\end{itemize}

% ============================================================
\section{Resultados: Todos los Estimadores}
% ============================================================

\subsection{Diseño A y B — Muestra completa}

\begin{threeparttable}
\begin{tabular}{l cc cc}
\toprule
 & \multicolumn{2}{c}{TIENE\_BILLETERA} & \multicolumn{2}{c}{USA\_BILLETERA} \\
\cmidrule(lr){2-3} \cmidrule(lr){4-5}
Especificación & Estimador & IC 95\% & Estimador & IC 95\% \\
\midrule
\multicolumn{5}{l}{\textit{Sharp RDD (ITT) --- Diseño A}} \\
MAIN\_baseline & $+0.0006$ & $[-0.036,\,+0.037]$ & $+0.0015$ & $[-0.034,\,+0.037]$ \\
\quad SE (robusto) & $(0.0187)$ & & $(0.0180)$ & \\
\quad N\_eff / BW & 35{,}006 / 18.04 & & 33{,}267 / 17.29 & \\[4pt]
MAIN\_covariates & $-0.0033$ & $[-0.031,\,+0.024]$ & $+0.0006$ & $[-0.020,\,+0.021]$ \\
\quad SE (robusto) & $(0.0139)$ & & $(0.0106)$ & \\
\quad N\_eff / BW & 29{,}139 / 15.85 & & 27{,}325 / 14.04 & \\[4pt]
MAIN\_dpto\_fe & $-0.0140$ & $[-0.027,\,-0.001]$ & $+0.0084$ & $[-0.001,\,+0.018]$ \\
\quad SE (robusto) & $(0.0065)$ & & $(0.0048)$ & \\
\quad N\_eff / BW & 57{,}039 / 33.46 & & 57{,}039 / 33.46 & \\
\midrule
\multicolumn{5}{l}{\textit{Fuzzy RDD (LATE) --- Diseño B}} \\
MAIN\_fuzzy & $-0.065$ & $[-0.645,\,+0.515]$ & $+0.143$ & $[-0.430,\,+0.717]$ \\
\quad SE (robusto) & $(0.296)$ & & $(0.293)$ & \\
\quad N\_eff / BW & 12{,}616 / 6.45 & & 12{,}616 / 6.07 & \\
\quad Método & rdrobust\_fuzzy & & rdrobust\_fuzzy & \\
\bottomrule
\end{tabular}
\begin{tablenotes}
\small
\item \textit{Notas:} Regresión local lineal, kernel triangular, bandwidth MSE-óptimo (rdrobust). SE robustos con sesgo corregido. MAIN\_dpto\_fe usa WLS local con efectos fijos de DPTO (fallback statsmodels por incompatibilidad de tipo en rdrobust con matrices dummy). Fuzzy RDD estima el LATE via Wald = (forma reducida) / (primera etapa) con tratamiento endógeno \texttt{RECIBE\_P65\_PERSONA}.
\end{tablenotes}
\end{threeparttable}

\begin{resultcaja}[Interpretación para la tesis — Diseño A (ITT)]
El ITT es prácticamente cero y estadísticamente no significativo en todas
las especificaciones del sharp RDD: el estimador oscila entre $-0.003$ y
$+0.001$ con SE en torno a 0.019. El intervalo de confianza al 95\% del
baseline es $[-0.036,\, +0.037]$, excluyendo efectos mayores a 3.7 puntos
porcentuales. Esto significa que la mera elegibilidad de edad para Pensión
65 no produce un aumento detectable en la adopción de billetera digital en
la población general peruana. El resultado es preciso, no solo no
significativo: la cota superior del IC excluye cualquier efecto de magnitud
práctica.
\end{resultcaja}

\begin{resultcaja}[Interpretación para la tesis — Diseño B (Fuzzy LATE)]
El LATE estimado es de $-0.065$ pp para TIENE\_BILLETERA y $+0.143$ pp para
USA\_BILLETERA, ambos con SE $\approx 0.30$ e IC que abarcan más de $\pm 0.60$.
Estos resultados no permiten concluir nada sobre el efecto causal de recibir
el bono: el rango de incertidumbre es demasiado amplio. La razón estructural
es que en el bandwidth estrecho ($h = 6.4$ años) la primera etapa es débil
---muy pocos adultos en el rango exacto 58--72 años reciben el bono---
lo que infla el denominador del Wald y la varianza del LATE. Nótese que el
signo negativo para TIENE\_BILLETERA no tiene interpretación causal confiable
dado el SE enorme.

\textbf{Para la tesis:} reportar el LATE como especificación alternativa con
la advertencia de instrumento débil en la full sample. El resultado primario
es el ITT del Diseño A.
\end{resultcaja}

\subsection{Diseño C — Muestra ELEGIBLE\_SIS=1 (IFH BCK 2017)}

\begin{center}
\begin{tabular}{lcccc}
\toprule
Especificación & Estimador & SE & IC 95\% & N\_eff / BW \\
\midrule
\multicolumn{5}{l}{\textit{TIENE\_BILLETERA}} \\
IFH\_sample\_baseline & $+0.019$ & 0.036 & $[-0.052,\,+0.090]$ & 2{,}218 / 9.92 \\
IFH\_sample\_covariates & $+0.013$ & 0.027 & $[-0.041,\,+0.066]$ & 2{,}185 / 9.00 \\
\multicolumn{5}{l}{\textit{USA\_BILLETERA}} \\
IFH\_sample\_baseline & $-0.021$ & 0.033 & $[-0.086,\,+0.044]$ & 1{,}802 / 7.20 \\
IFH\_sample\_covariates & $-0.024$ & 0.023 & $[-0.070,\,+0.021]$ & 1{,}553 / 6.22 \\
\bottomrule
\end{tabular}
\end{center}

\begin{interpcaja}[Interpretación para la tesis]
Los estimadores del Diseño C son positivos para TIENE\_BILLETERA (+0.019,
+0.013) pero estadísticamente no significativos: los IC cruzan cero
ampliamente. La dirección del efecto es la esperada (el programa aumenta
la adopción entre los elegibles SISFOH), pero el tamaño muestral restringido
($N_{\text{eff}} \approx 2{,}200$ dentro del BW) no da suficiente potencia
para rechazar el nulo. Reportar como análisis de heterogeneidad sugestivo.
\end{interpcaja}

\subsection{Heterogeneidad descriptiva: extrema pobreza monetaria (POBREZA=1)}

\begin{center}
\begin{tabular}{lcccc}
\toprule
Outcome & Estimador & SE & IC 95\% & N\_eff / BW \\
\midrule
TIENE\_BILLETERA & $+0.028$ & 0.021 & $[-0.014,\,+0.070]$ & 1{,}144 / 13.79 \\
USA\_BILLETERA & $-0.002$ & 0.001 & $[-0.003,\,-0.000]$ & 663 / 7.62 \\
\bottomrule
\end{tabular}
\end{center}

\begin{alertcaja}[Advertencia para la tesis]
Estos son resultados \textbf{descriptivos, no causales}. La muestra POBREZA=1
excluye al 88.7\% de los receptores reales del bono (los receptores tienen
mayor gasto por el bono mismo, lo que los saca de la categoría
POBREZA=1). Cualquier discontinuidad estimada en esta submuestra mezcla
el efecto del programa con el sesgo de selección ex-post. Reportar con
esta advertencia explícita.
\end{alertcaja}

\subsection{Heterogeneidad urbano/rural}

\begin{center}
\begin{tabular}{lccccc}
\toprule
Subgrupo & N total & Estimador & SE & IC 95\% & N\_eff / BW \\
\midrule
Urbano (AREA=1) & 41{,}568 & $-0.032$ & 0.036 & $[-0.103,\,+0.038]$ & 8{,}800 / 12.08 \\
Rural (AREA=2) & 58{,}342 & $+0.016$ & 0.020 & $[-0.022,\,+0.054]$ & 11{,}211 / 11.79 \\
\bottomrule
\end{tabular}
\end{center}

\begin{interpcaja}[Interpretación para la tesis]
Existe una diferencia de signo entre áreas: el efecto es ligeramente
negativo en urbano ($-0.032$) y positivo en rural ($+0.016$), aunque ambos
son estadísticamente nulos. La dirección del efecto en rural es plausible:
en zonas rurales la billetera digital puede ser la única forma de cobrar el
bono (Cuenta DNI), creando un efecto de adopción inducida. En urbano, el
efecto puede ser negativo por composición etaria (adultos mayores urbanos
tienen mayor probabilidad previa de usar efectivo). Sin embargo, estos
efectos son demasiado imprecisos para conclusiones causales.
\end{interpcaja}

\subsection{Diseño D — RDD2: IFH como running variable}

\begin{center}
\begin{tabular}{lccccc}
\toprule
Especificación & Outcome & Estimador & SE & IC 95\% & N\_eff / BW \\
\midrule
RDD2\_IFH\_baseline & TIENE\_BILLETERA & $+0.028$ & 0.021 & $[-0.013,\,+0.069]$ & 5{,}052 / 12.48 \\
RDD2\_IFH\_covariates & TIENE\_BILLETERA & $+0.027$ & 0.018 & $[-0.009,\,+0.063]$ & 4{,}653 / 11.52 \\
RDD2\_IFH\_baseline & USA\_BILLETERA & $+0.015$ & 0.024 & $[-0.031,\,+0.061]$ & 4{,}479 / 10.76 \\
RDD2\_IFH\_covariates & USA\_BILLETERA & $+0.011$ & 0.015 & $[-0.019,\,+0.040]$ & 4{,}505 / 11.03 \\
\midrule
\multicolumn{6}{l}{\textit{Balance de covariables en umbral IFH (adultos $\geq$65)}} \\
INTERNET\_HOGAR & --- & $+0.055$ & 0.155 & $[-0.249,\,+0.360]$ & 6{,}178 / 15.71 \\
NIVEL\_EDUCATIVO & --- & $+0.047$ & 0.976 & $[-1.867,\,+1.961]$ & 8{,}190 / 23.70 \\
INGRESO\_PC & --- & $+22.8$ & 3{,}685 & $[-7{,}200,\,+7{,}246]$ & 8{,}028 / 22.26 \\
\bottomrule
\end{tabular}
\end{center}

\begin{resultcaja}[Interpretación para la tesis — Diseño D (RDD2)]
El RDD2 sugiere un efecto positivo de $\approx +2.8$ pp sobre TIENE\_BILLETERA
para adultos mayores clasificados como elegibles SISFOH (IFH justo bajo el
umbral) versus no elegibles (IFH justo sobre el umbral). Sin embargo, el IC
al 95\% incluye el cero (baseline: $[-0.013,\, +0.069]$), por lo que no se
rechaza el nulo estadístico.

Importante: el balance de covariables es satisfactorio en el umbral IFH
(todos los test de balance arrojan estimadores NS). Esto valida la
continuidad local del IFH como running variable.

El RDD2 es el diseño más cercano a la estrategia de Bando, Galiani y
Gertler (2020) para Perú, y la dirección del efecto (+2.8 pp) es
económicamente relevante aunque imprecisa. Nótese que la cota superior
del IC es $+0.069$ (6.9 pp), lo que no descarta efectos de magnitud práctica.
\end{resultcaja}

% ============================================================
\section{Pruebas de Robustez}
% ============================================================

\subsection{Panel A: Sensibilidad al bandwidth}

\begin{center}
\begin{tabular}{lrcrrl}
\toprule
Especificación & Estimador & SE & IC 95\% & BW & N\_eff \\
\midrule
$h = h^*/2$ (mitad óptimo) & $-0.006$ & 0.025 & $[-0.054,\,+0.043]$ & 9.02 & 18{,}649 \\
$h = h^*$ (óptimo, baseline) & $-0.002$ & 0.019 & $[-0.040,\,+0.036]$ & 18.04 & 35{,}006 \\
$h = 2h^*$ (doble óptimo) & $+0.006$ & 0.016 & $[-0.026,\,+0.038]$ & 36.09 & 62{,}536 \\
\bottomrule
\end{tabular}
\end{center}

\begin{interpcaja}[Interpretación]
El estimador es estable alrededor de cero para todo el rango de BW
evaluado: entre $-0.006$ y $+0.006$. Los IC siempre incluyen el cero.
La estabilidad del nulo ante variaciones del BW es una evidencia fuerte
de que el resultado no es un artefacto de la selección del BW.
\end{interpcaja}

\subsection{Panel B: Donut-hole (exclusión de observaciones en el cutoff)}

\begin{center}
\begin{tabular}{lrcrrl}
\toprule
Exclusión & Estimador & SE & IC 95\% & BW & N\_eff \\
\midrule
$|\text{EDAD} - 65| \leq 0.5$ & $+0.008$ & 0.021 & $[-0.034,\,+0.049]$ & 18.04 & 33{,}947 \\
$|\text{EDAD} - 65| \leq 1.0$ & $-0.009$ & 0.021 & $[-0.052,\,+0.033]$ & 18.04 & 31{,}941 \\
$|\text{EDAD} - 65| \leq 2.0$ & $+0.008$ & 0.030 & $[-0.050,\,+0.065]$ & 18.04 & 30{,}026 \\
\bottomrule
\end{tabular}
\end{center}

\begin{interpcaja}[Interpretación]
La especificación donut elimina observaciones en un entorno del cutoff para
mitigar problemas de manipulación local. Los resultados son nulos y
estables, descartando que el efecto sea impulsado por observaciones con
exactamente 65 años de edad.
\end{interpcaja}

\subsection{Panel C: Sensibilidad al grado del polinomio}

\begin{center}
\begin{tabular}{lrcrrl}
\toprule
Polinomio & Estimador & IC 95\% & BW & N\_eff \\
\midrule
Lineal (baseline) & $+0.001$ & $[-0.036,\,+0.037]$ & 18.04 & 35{,}006 \\
Cuadrático & $-0.007$ & $[-0.050,\,+0.036]$ & 16.51 & 31{,}464 \\
\bottomrule
\end{tabular}
\end{center}

\subsection{Panel D: Placebos de covariables como outcomes}

\begin{center}
\begin{tabular}{lrcrrl}
\toprule
Covariable (como outcome) & Estimador & SE & IC 95\% & BW & N\_eff \\
\midrule
INTERNET\_HOGAR & $+0.034$ & 0.049 & $[-0.062,\,+0.130]$ & 15.28 & 29{,}668 \\
SMARTPHONE & $+0.027$ & 0.025 & $[-0.022,\,+0.075]$ & 14.30 & 27{,}809 \\
\bottomrule
\end{tabular}
\end{center}

\begin{interpcaja}[Interpretación]
Las covariables preexistentes (internet y smartphone) no muestran
discontinuidades significativas en el cutoff de 65 años. Esto confirma
que las características del hogar son continuas en el umbral, validando el
supuesto de continuidad del RDD.
\end{interpcaja}

\subsection{Panel E: Balance de covariables en el bandwidth}

\begin{center}
\begin{tabular}{lrcrrl}
\toprule
Covariable & Estimador & SE & IC 95\% & BW & N\_eff \\
\midrule
INTERNET\_HOGAR & $+0.029$ & 0.048 & $[-0.065,\,+0.123]$ & 18.04 & 35{,}006 \\
SMARTPHONE & $+0.029$ & 0.025 & $[-0.020,\,+0.078]$ & 18.04 & 35{,}006 \\
INGRESO\_PC & $+362$ & 1{,}314 & $[-2{,}214,\,+2{,}938]$ & 18.04 & 35{,}006 \\
NIVEL\_EDUCATIVO & $+0.120$ & 0.304 & $[-0.477,\,+0.716]$ & 18.04 & 34{,}392 \\
\bottomrule
\end{tabular}
\end{center}

\begin{resultcaja}[Interpretación]
Ninguna covariable muestra un salto estadísticamente significativo en el
cutoff de 65 años. Este resultado valida el supuesto clave del RDD: que
las características preexistentes de los individuos son continuas en la
running variable. Un diseño de RDD válido requiere que solo el outcome
(y el tratamiento, en fuzzy) salten en el cutoff.
\end{resultcaja}

\subsection{Panel F: Placebo de cutoff}

\begin{center}
\begin{tabular}{lrcrl}
\toprule
Cutoff placebo & Estimador & IC 95\% & BW & Resultado \\
\midrule
$c = -50$ (izquierda) & $-0.007$ & $[-0.007,\,-0.006]$ & 3.09 & Rechaza nulo* \\
$c = -12$ (derecha) & $+0.019$ & $[-0.020,\,+0.058]$ & 6.09 & No rechaza \checkmark \\
\bottomrule
\end{tabular}
\end{center}

\textbf{Nota:} el resultado del cutoff izquierdo (c = $-50$, equivalente a
edad 15) puede interpretarse como un artifact estadístico por el BW muy
estrecho (3.09 años), no como evidencia de manipulación.

\subsection{Panel G: Inferencia por permutación}

\textbf{Estadístico real:} $t = -2.84$ \\
\textbf{p-valor permutacional} (500 iteraciones): \textbf{0.94}

\begin{interpcaja}[Interpretación]
El p-valor permutacional de 0.94 indica que el estadístico observado
($t = -2.84$) está dentro del rango de variación que se obtendría si el
cutoff fuera aleatorio. No se rechaza la hipótesis nula. Este es uno de
los tests más robustos del análisis, ya que no depende de supuestos de
distribución.
\end{interpcaja}

% ============================================================
\section{Figuras Generadas}
% ============================================================

\begin{longtable}{@{}llp{8cm}@{}}
\toprule
Archivo & Sección del paper & Descripción \\\midrule
\endhead
\texttt{figure\_1\_rdplot.pdf} & Figura 1 & RD plot principal: media de TIENE\_BILLETERA por bin de edad, con la discontinuidad estimada en los 65 años. Puntos = medias de bins; línea = ajuste polinomial local lineal a cada lado del cutoff. \\[6pt]
\texttt{figure\_2\_bandwidth\_sensitivity.pdf} & Figura 2 & Estimadores del ITT para un rango de bandwidths $h \in [7, 36]$ con IC al 95\%. Muestra estabilidad del nulo. \\[6pt]
\texttt{figure\_bandwidth\_sensitivity\_granular.pdf} & Apéndice & Versión más granular de la sensibilidad al BW. \\[6pt]
\texttt{figure\_mccrary\_density.pdf} & Figura 3 & Test de densidad McCrary/rddensity: histograma de la densidad de EDAD alrededor del cutoff. Permite detectar manipulación de la running variable. \\[6pt]
\texttt{figure\_first\_stage.pdf} & Figura 4 & RD plot de la primera etapa: RECIBE\_P65\_PERSONA vs EDAD. Muestra el salto en recepción del bono en los 65 años. \\[6pt]
\texttt{figure\_descriptivos\_brecha.pdf} & Apéndice & Brechas descriptivas en adopción digital por grupos de edad, urbano/rural, POBREZA. \\[6pt]
\texttt{figure\_pobreza\_receptores.pdf} & Apéndice & Distribución de receptores de Pensión 65 por categoría POBREZA. Visualiza el 88.7\% excluido. \\[6pt]
\texttt{figure\_rdd2\_rdplot.pdf} & Figura 5 & RD plot del Diseño D (RDD2): TIENE\_BILLETERA vs IFH centrado en umbral. Ventana $\pm 30$ puntos IFH. \\[6pt]
\texttt{figure\_rdd2\_density.pdf} & Figura 6 & Test McCrary para el RDD2: densidad del IFH en el umbral departamental. \\
\bottomrule
\caption{Figuras generadas por el pipeline. Todas en \texttt{paper/figures/}.}
\end{longtable}

% ============================================================
\section{Tablas LaTeX Generadas}
% ============================================================

\begin{center}
\begin{tabular}{llp{7cm}}
\toprule
Archivo & Tabla en paper & Contenido \\\midrule
\texttt{table\_1\_summary.tex} & Tabla 1 & Estadísticas descriptivas por grupo tratado/control \\
\texttt{table\_2\_main\_results.tex} & Tabla 2 & Resultados principales: MAIN\_baseline, MAIN\_covariates, MAIN\_fuzzy, MAIN\_dpto\_fe \\
\texttt{table\_3\_heterogeneity.tex} & Tabla 3 & Heterogeneidad: POBREZA=1, urbano, rural, IFH sample \\
\texttt{table\_3\_robustness.tex} & Tabla 3b & Robustez: BW, donut, polinomios, placebos \\
\texttt{table\_4\_covariate\_balance.tex} & Tabla 4 & Balance de covariables en cutoff \\
\texttt{table\_4\_rdd2\_ifh.tex} & Tabla 5 & RDD2 con IFH como running variable (BCK 2017) \\
\texttt{table\_A1\_ep\_descriptivo.tex} & Apéndice A1 & Resultados descriptivos POBREZA=1 con advertencia \\
\bottomrule
\end{tabular}
\end{center}

% ============================================================
\section{Síntesis de Hallazgos para la Tesis}
% ============================================================

\subsection{Resultado principal}

\textbf{El efecto causal de la elegibilidad de Pensión 65 sobre la adopción
de billetera digital en la población general peruana (ENAHO 2024) es
estadísticamente nulo.}

El estimador ITT puntual es $\hat{\tau} = +0.0006$ pp (SE $= 0.019$,
IC 95\% $= [-0.036,\, +0.037]$). Este resultado es preciso: la cota
superior del IC excluye cualquier efecto mayor a 3.7 pp.

\subsection{Resultados secundarios}

\begin{enumerate}
  \item \textbf{Primera etapa válida:} la edad de 65 años genera un salto
    nítido en la recepción del bono (de 0\% a 3.8\% exactamente en el
    cutoff, subiendo a 16\% a los 66 años). El instrumento existe, pero la
    dilución en la full sample es grande.
  \item \textbf{RDD2 con IFH:} para adultos mayores elegibles según el
    puntaje IFH de Bernal et al. (2017), el efecto sobre TIENE\_BILLETERA
    es $+0.028$ pp (IC $[-0.013,\, +0.069]$). No se rechaza el nulo pero
    la cota superior es 6.9 pp, lo que no descarta efectos de magnitud
    práctica.
  \item \textbf{Heterogeneidad por área:} en zonas rurales el efecto es
    positivo ($+0.016$ pp) mientras que en urbanas es negativo ($-0.032$ pp),
    aunque ambos son no significativos. La dirección rural es plausible dado
    que la Cuenta DNI es el mecanismo de cobro en áreas sin agencias bancarias.
  \item \textbf{LATE impreciso:} el Fuzzy RDD con RECIBE\_P65\_PERSONA estima
    un LATE de $-0.065$ pp pero con SE $= 0.296$ (problema de instrumento
    débil en el BW estrecho). No es informativo sobre el efecto causal.
  \item \textbf{Robustez del nulo:} 18 de los 21 tests de robustez confirman
    el resultado nulo (NS). El test de permutación da p $= 0.94$.
\end{enumerate}

\subsection{Interpretación causal}

El nulo preciso tiene tres interpretaciones no mutuamente excluyentes:

\begin{enumerate}
  \item \textbf{Efecto real nulo:} recibir Pensión 65 no aumenta la adopción
    de billetera digital. Los beneficiarios cobran el bono en ventanillas
    físicas de los bancos y no adoptan la banca digital.
  \item \textbf{Poder estadístico insuficiente para detectar efectos pequeños:}
    dado el bajo take-up en el bandwidth ($\approx 5\%$ de la muestra completa
    tiene 65 años \textit{y} recibe el bono), el diseño sharp RDD sobre la
    full sample está underpowered para detectar efectos menores a 3.7 pp.
  \item \textbf{Confound de cohorte:} en datos cross-section, comparar personas
    de 64 vs 66 años equivale a comparar cohortes nacidas en 1960 vs 1958,
    con diferente exposición de vida a la tecnología digital. Este confound
    no puede eliminarse con el diseño actual sin datos de panel o una
    encuesta antes del despliegue de Yape/Plin (2018--2020).
\end{enumerate}

\subsection{Limitaciones a documentar en la tesis}

\begin{enumerate}[noitemsep]
  \item La running variable (EDAD) es discreta (enteros), lo que genera
    ``mass points''. rdrobust lo detecta y lo maneja, pero puede reducir
    la precisión del BW MSE-óptimo.
  \item POBREZA $\neq$ SISFOH: el análisis no puede replicar exactamente el
    criterio de elegibilidad SISFOH porque INEI no libera el puntaje oficial.
    El IFH de BCK (2017) es una aproximación con 15 clusters; se usan
    25 departamentos.
  \item El diseño cross-section no permite identificar el efecto de corto
    vs. largo plazo de la adopción inducida por el bono.
  \item Discontinuidades competidoras potenciales a los 65 años: AFP, ONP,
    SIS$+$65, pensiones privadas. No se probaron explícitamente (excepción:
    el placebo de Juntos pasa con NS).
\end{enumerate}

% ============================================================
\section{Especificaciones Técnicas del Estimador}
% ============================================================

\begin{center}
\begin{tabular}{ll}
\toprule
Parámetro & Valor \\\midrule
Software & Python 3.14, \texttt{rdrobust 0.x}, \texttt{statsmodels}, \texttt{pandas} \\
Kernel & Triangular \\
Selección de BW & MSE-óptimo (\texttt{bwselect = ``mserd''}) \\
Errores estándar & HC1 robustos con sesgo corregido (rdrobust) \\
Polinomio base & Local lineal ($p=1$) \\
Semilla aleatoria & 42 \\
Clustering & \texttt{CONGLOME} (conglomerado muestral) \\
BW baseline (TIENE\_BILLETERA) & $h^* = 18.04$ años \\
BW baseline (USA\_BILLETERA) & $h^* = 17.29$ años \\
\bottomrule
\end{tabular}
\end{center}

\bigskip
\noindent\textbf{Archivos clave del repositorio:}
\begin{itemize}[noitemsep]
  \item \texttt{scripts/script\_completo/script.py} --- pipeline principal
  \item \texttt{construir\_ifh.py} --- construcción del IFH (BCK 2017)
  \item \texttt{data/clean/main\_results.csv} --- 18 estimaciones principales
  \item \texttt{data/clean/robustness\_results.csv} --- 21 tests de robustez
  \item \texttt{data/clean/rdd2\_results.csv} --- 7 estimaciones RDD2
  \item \texttt{paper/tables/} --- todas las tablas LaTeX listas para insertar
  \item \texttt{paper/figures/} --- todas las figuras en PDF y PNG
\end{itemize}

\bigskip
\hrule
\bigskip
{\small\textit{Documento generado automáticamente el 2026-07-02 a partir de los
outputs verificados del pipeline. Todas las cifras se leen directamente de
\texttt{main\_results.csv}, \texttt{robustness\_results.csv} y
\texttt{rdd2\_results.csv}. Autor del análisis: Cecilia Vargas Risco.}}

\end{document}
